\section{Compatible joint certificate and fault-contingent semantics}
\label{sec:joint-rigorous}
For each proposed switching action \(a\), fix a nonempty uncertainty set \(\Omega_a\), an observation horizon \(H_a>0\), and a finite set \(\mathcal F_a\) of relevant faults. Each scenario \(\omega\in\Omega_a\) specifies a frequency input on \([0,H_a]\), and, for each fault \(\ell\), a counterfactual sustained-fault relay experiment with a primary \(p_\ell\), backup \(b_\ell\), relay model, and initial timer state. The relay model must be fixed for each pair: constant-current inverse-time or time-domain definite-time with complete reset; a numerical operating-time inequality from the former does not prove pickup persistence in the latter.

Let \(t_{p,\ell}(\omega)\) and \(t_{b,\ell}(\omega)\) denote finite relay assertion times under \emph{the same sustained-fault scenario}; let \(d_{p,\ell}(\omega)\) and \(d_{b,\ell}(\omega)\) denote nonnegative breaker opening delays. Define total clearing times \(c_{r,\ell}=t_{r,\ell}+d_{r,\ell}\), for \(r\in\{p,b\}\). The backup time is counterfactual: after successful primary isolation, a backup relay should be inhibited or should not reach its trip condition. Let \(D_{p,\ell}\), \(D_{b,\ell}\) be the prescribed clearing deadlines and \(\Delta_\ell>0\) the coordination margin. Assume that the primary and backup operate on the specified fault currents and that relay assertion and breaker opening are correctly represented by the selected model.

\begin{theorem}[Joint model-relative sufficiency]\label{thm:joint-suff}
Suppose for a given action \(a\): (i) its scalar frequency dynamics satisfy the assumptions of Proposition~\ref{prop:frequency} for every \(\omega\in\Omega_a\), and the finite-horizon frequency envelope does not exceed \(f_{\max}\); (ii) for each \(\ell\in\mathcal F_a\) and each \(\omega\in\Omega_a\), the stipulated relay model gives finite times satisfying
\[
c_{p,\ell}(\omega)\le D_{p,\ell},\qquad
c_{b,\ell}(\omega)\le D_{b,\ell},\qquad
t_{b,\ell}(\omega)-c_{p,\ell}(\omega)\ge\Delta_\ell;
\]
and (iii) whenever the selected relay model has a pickup-persistence or reset requirement, the required operating times are certified using that \emph{same} time-domain model and initial state. Then the action satisfies the frequency envelope and the specified primary clearing, hypothetical backup clearing and non-premature-backup coordination requirements for all \(\omega\in\Omega_a\). This theorem makes no assertion about voltage recovery, switch transients, fault detection, communications, or real devices outside the model.
\end{theorem}
\begin{proof}
Fix an arbitrary \(\omega\in\Omega_a\). Proposition~\ref{prop:frequency} bounds its frequency trajectory throughout the horizon. For every fault \(\ell\), condition (ii) directly bounds both clearing times and ensures that backup relay assertion occurs at least \(\Delta_\ell\) after primary breaker opening in the sustained-fault counterfactual. Condition (iii) guarantees that the operating times used in (ii) are generated by the chosen relay semantics rather than by an incompatible static approximation. The scenario and fault were arbitrary, so all asserted model constraints hold universally.
\end{proof}
\paragraph{Operational separation.}
The theorem is a \emph{sufficiency composition result}, not a new physical plant theorem. In particular, the coordination margin compares backup \emph{assertion} against primary \emph{clearing}; comparing two relay assertion times alone omits primary breaker opening. If the backup is armed only after a detection or communications delay, that mechanism must enter its operating-time model. Independent uncertainty intervals are safe outer approximations even if physically correlated, but they may be conservative. To use an interval extremum, breaker delays must be bounded consistently: worst primary clearing uses its largest breaker delay; earliest backup assertion uses its smallest allowed delay. A finite backup clearing deadline is checked separately with its largest delay.

\begin{proposition}[Full-reset pickup equivalence]\label{prop:reset}
Let \(I(t)\) be piecewise constant on a finite partition of \([0,D]\), with zero initial timer, strict pickup \(I(t)>I_{\rm pk}\), positive required dwell \(T_{\rm req}\), instantaneous full reset whenever pickup is lost, and an inclusive deadline \(D\). The relay asserts by \(D\) if and only if there exists a contiguous above-pickup interval of length at least \(T_{\rm req}\) contained in \([0,D]\).
\end{proposition}
\begin{proof}
The timer increases at unit rate only while pickup holds, and otherwise resets to zero. Hence it reaches \(T_{\rm req}\) precisely when an uninterrupted above-pickup interval has accumulated that duration. Neither separate shorter intervals nor an isolated current peak can trigger the stipulated relay.
\end{proof}
